\chapter{实验验证与分析}

\section{实验设计与数据集}

\subsection{数据采集}\label{sec:data_collection}

实验数据通过第三章所述的无人机多模态实验平台采集。数据采集在室内受控环境中进行，被测表面为铝板（厚度3 mm，平面尺寸$600 \times 600$ mm）。单次飞行实验包含完整的接近-接触-保持-脱离流程，由位置-速度串级PID控制器驱动无人机从距表面约1.5 m处出发，以~0.2 m/s速度向目标表面接近。

数据采集系统在每次飞行中同步记录以下数据：（1）EMAT波形数据（~40 Hz，~32740采样点/帧，1 MHz采样率）；（2）FAST-LIO里程计数据（~100 Hz，6-DOF位姿+速度）；（3）D435彩色图像和对齐深度图（~30 Hz，640$\times$480分辨率）。共计采集[N]次有效飞行实验，累计[N]分钟多模态数据。% 待填充具体数据

\subsection{标签生成}

无人机接触检测不存在天然完美标签。为此，本研究采用弱标签与人工复核结合的策略构建训练标签。初始标签由以下规则联合生成：（1）当位置误差$\|\mathbf{e}_p\|_2 < 0.05$ m且法向速度$|v_n| < 0.02$ m/s时，标记为“接触候选”状态；（2）当EMAT回波能量超过静息态3倍标准差且频谱重心稳定时，确认“信号有效接触”；（3）视觉深度图上目标ROI区域的深度值变化率小于5\%/秒。初始标签再经人工逐帧复核，重点修正接近-接触过渡段和接触-脱离过渡段的边界标签，使模型学习到连续过程而非孤立的帧标签。

\section{接触判别性能评估}

\subsection{评价指标}

采用以下五项指标评估接触判别性能：

\begin{itemize}
\item \textbf{逐帧准确率}（Accuracy）：所有时间帧中正确判别的比例。
\item \textbf{接触状态F1}（F1-score）：针对接触类别的精确率与召回率的调和均值，更侧重反映模型在接触段的表现。
\item \textbf{临界状态误报率}（Critical FPR）：在接触建立和脱离前后$\pm 2$秒的临界区间内的误报率，反映模型在过渡段的判别可靠性。
\item \textbf{时序平滑度}（Temporal Smoothness）：相邻帧接触概率的均方根变化量$\sqrt{\frac{1}{T-1}\sum_t (\hat{p}_{t+1} - \hat{p}_t)^2}$，值越小表明概率序列越平滑。
\item \textbf{单帧推理延迟}（Inference Latency）：从多模态对齐特征输入到接触概率输出的单帧耗时（ms），反映算法的实时性。
\end{itemize}

\subsection{与基线方法的比较}\label{sec:baseline_comparison}

将本文提出的完整模型（PhysConAttn + Gating + Smooth）与以下基线方法进行比较：

\begin{itemize}
\item \textbf{单模态超声阈值法}：仅使用EMAT回波能量的固定阈值（>3$\sigma$静息态）进行接触判别，代表传统方法。
\item \textbf{特征拼接+MLP}：将三模态特征直接拼接后输入三层MLP分类器，不包含时序建模和物理约束。
\item \textbf{标准Transformer（无物理约束）}：使用标准自注意力Transformer进行时序建模，但不注入物理约束矩阵。
\item \textbf{Transformer+物理约束（无门控）}：注入物理约束注意力但使用固定等权融合，不包含模态可靠性门控。
\item \textbf{本文完整模型}：物理约束注意力 + 模态可靠性门控 + 时序平滑损失。
\end{itemize}

实验结果以表格形式汇总如下。% 待填充具体性能数据

\begin{table}[htbp]
\centering
\caption{各方法接触判别性能对比}
\label{tab:results}
\begin{tabular}{lccccc}
\toprule
\textbf{方法} & \textbf{Acc.(\%)} & \textbf{F1} & \textbf{Crit.FPR} & \textbf{Smooth.} & \textbf{Lat.(ms)} \\
\midrule
超声单阈值 & -- & -- & -- & -- & -- \\
特征拼接+MLP & -- & -- & -- & -- & -- \\
标准Transformer & -- & -- & -- & -- & -- \\
Transformer+物理约束 & -- & -- & -- & -- & -- \\
\textbf{本文完整模型} & \textbf{--} & \textbf{--} & \textbf{--} & \textbf{--} & \textbf{--} \\
\bottomrule
\end{tabular}
\end{table}

\textit{注：表中数据待实验完成后填充。}

\section{消融实验}

按照第五章设计的消融方案，逐一去除各模块并进行性能对比。预期结果如下：

\begin{enumerate}
\item \textbf{去除物理约束矩阵}：预期在临界接触段的误报率显著升高，注意力权重分布趋向均匀化或集中于噪声帧，失去时间邻近性和空间邻近性的引导。
\item \textbf{去除模态可靠性门控}：预期在模态退化场景（如视觉反光、姿态未收敛）下性能下降更为明显，门控权重的可视化应显示各模态权重随退化程度单调变化。
\item \textbf{去除时序平滑损失}：预期相邻帧接触概率出现高频闪烁，即使在总体准确率变化不大的情况下，时序平滑度指标应明显恶化。
\end{enumerate}

\section{无人机闭环扫查验证}

为验证多模态接触感知在真实无人机飞行中的有效性，在铝板上方进行自主接近-保持实验。无人机从距表面约1.5 m处自主接近，当本文模型的接触概率连续3秒（约120帧）超过$\tau_{\text{on}} = 0.85$时，飞控状态机从APPROACH切换至HOLD，保持当前位置；当概率持续低于$\tau_{\text{off}} = 0.4$且伴随法向残差增大时，状态机回到APPROACH。

关键观测指标包括：（1）状态切换延迟——从首次满足接触条件到状态机实际切换的时间差；（2）闭环扫查稳定性——在HOLD状态下法向位置的标准差；（3）状态闪烁次数——在临界接触段内单帧状态翻转的次数。

\section{本章小结}

本章设计了系统性的实验验证方案，并给出了初步的实验环境设置和评估框架。从单帧分类准确率到闭环飞行稳定性，评价指标构成了层次化的性能评估体系。消融实验设计确保每个模块的功能贡献具有可证伪性。需要说明的是，本章中标注为“--”的数值结果待后续实验完成后补充。

% ══════════════════════════════════════════════════════════════════
% 第七章 总结与展望
% ══════════════════════════════════════════════════════════════════